\subsection{Deep Reinforcement Learning Implementation}

\subsubsection{Deep Q-Network (DQN) Architecture}

The agent's decision-making policy is approximated using a Deep Q-Network
(DQN)~\cite{mnih2015humanlevel}. We employ a Multi-Layer Perceptron (MLP) architecture to
map the high-dimensional state vector
\(\mathbf{S}_t\) to Q-values for the five discrete actions.

The network is a 3-layer MLP with hidden dimensions [512, 512, 256] and ReLU activations.
The input dimension is $4(3+2N)$ (frame-stacked state) and the output is 5 Q-values,
one per action in $\mathcal{A}$.

The network weights \(\theta\) are optimised using the Adam optimiser to minimise the temporal
difference error between the predicted Q-values and the target Q-values derived from the
Bellman equation. To ensure training stability, a separate target network \(\hat{Q}\) with
parameters \(\theta^{-}\) is maintained and updated periodically (every \(C\) steps).

\subsubsection{Probabilistic Link-Layer Model}

Rather than binary ``in-range'' abstractions, the environment models each sensor's transmission
success as the product of its protocol-mandated duty cycle and a sigmoid-mapped SNR.
The instantaneous success probability $P_{\text{overall}}$ is computed by comparing the RSSI
(derived from the Two-Ray Ground Reflection model) against the noise
floor and the SF-specific SNR threshold. Setting a threshold of $P_{\text{overall}} > 0.5$
forces the DQN agent to navigate the trade-off between proximity-induced signal strength and
the stochastic availability of sensor nodes, enabling active management of link quality
through positioning.

\subsubsection{Training Convergence}

The DQN was trained across three curriculum stages progressively introducing all 16 training
conditions (4 grid sizes $\times$ 4 sensor counts). Within each stage, the mean episode
reward follows a characteristic rise-and-plateau trajectory: early episodes produce large
negative rewards as the agent explores randomly, followed by a monotonic increase as the
policy learns to collect data, converging to a stable plateau. Later stages converge to
progressively higher absolute rewards ($3$--$5\times10^6$), confirming that the curriculum
successfully bootstraps the agent from easy to hard conditions.

The TD loss exhibits a spike-then-converge profile at each stage transition---the spike
reflects the temporary distribution shift as new grid sizes and sensor counts enter the
training mix, while convergence confirms the replay buffer and target network stabilise
learning despite this non-stationarity. Critically, no divergence occurs across any stage
transition, validating the curriculum ordering (small grids with few sensors first,
progressing to large grids with dense deployments). Full convergence curves for both reward
and TD loss across all stages are provided in Appendix~\ref{app:training}
(Figure~\ref{fig:training_convergence}).
